\chapter{Check-up and characterisation (menu 1)}
\label{ch:menu1}
\menulabel{1}

\section{Purpose}

Turn one ORCA output file into an analysis + diagnosis report: orbital
composition, occupation and entropy analysis, a multi-reference character
panel, the local spin analysis when the input defines fragments, and the
diagnostic findings with their provenance. This is the entry point for ``I have
an output file -- tell me what I actually got''.

\section{What you need}

An ORCA 6.x output file (\file{.out} or \file{.log}). Which sections appear
depends on what the output contains:

\begin{itemize}
  \item the \textbf{per-MO composition} analysis needs the
    \code{LOEWDIN ORBITAL-COMPOSITIONS} table: add
    \code{\%output Print[P\_ReducedOrbPopMO\_L] 1} to the input (it is printed
    at the normal print level); without the table the section is simply
    skipped, with an explanation in the report;
  \item the \textbf{occupation/entropy} analyses need a CASSCF output (the
    \code{N(occ)=} line);
  \item the \textbf{diffuse-orbital check} (A8) additionally needs the basis
    exponents themselves: add \code{!PrintBasis} to the input, so the run prints
    its shells in input format; without it the section is skipped (the exponents
    cannot be inferred from the composition table);
  \item the \textbf{T1} line of the multi-reference panel needs a coupled-cluster
    output;
  \item the \textbf{local spin} section appears only when the input divided the
    molecule into fragments (\code{N(1)} style atom labels, see
    Chapter~\ref{ch:menu1}, Section~\ref{sec:menu1-local-spin}).
\end{itemize}

Nothing else is required; a plain DFT output gives a report with the findings
that apply to it and nothing else.

\section{Walkthrough}

\lstinputlisting[style=fbkcode,
  title={The menu-1 example session (\file{examples/scripts/m1\_checkup.txt})}]
  {examples/expected/m1_checkup.txt}

The session writes two files beside the input, \file{name.fbk.md} and
\file{name.fbk.json}, and prints the section counts and a finding summary. The
body of the report for this example is walked through line by line in
Chapter~\ref{ch:report}; the following sections point at the variants.

\section{What you get}

\begin{itemize}
  \item \file{name.fbk.md} -- the report: summary, one section per applicable
    analyser (A1 orbital composition, A2 occupation/entropy bound, A3
    multi-reference panel, A6 local spin analysis, A8 diffuse-orbital check),
    findings ordered by severity, the provenance list, and the references block
    with complete citations and paste-ready BibTeX entries;
  \item \file{name.fbk.json} -- the same content as data, for scripting.
\end{itemize}

\subsection*{The Gaussian leg (the minimal fact set)}

Gaussian 09/16 outputs are identified by their ``Entering Gaussian System''
banner and mapped onto the shared diagnosis vocabulary -- termination, SCF
(``SCF Done'' shape), frequency blocks (the ``Frequencies --'' lines; the
engine's ``N imaginary frequencies'' mark is kept alongside the counted
value) and the optimization cycle facts (``Step number N out of a maximum
of M'', ``Optimization completed.''/``Stationary point found.'', and the
``ts'' token of the echoed input line).  Engine-independent rules then fire
normally -- the example's TS optimisation without a frequency verification
is caught by the same rule as on ORCA files -- while ORCA-only analyses and
rules stay absent.  Boundaries: Gaussian CASSCF facts are registered as a
follow-up (no sample in hand); the ``SCF Done'' parse follows the published
line shape with its real fixture anchor registered (the two shipped G09
probes are external-driver runs without an SCF line); the suggested-action
text of shared rules may quote ORCA syntax -- the rule and its reasoning
are engine-independent, the syntax is not.

\lstinputlisting[style=fbkcode, firstline=146, lastline=151,
  title={The Gaussian leg of the menu-1 example (a G09 TS optimisation
  without frequency verification)}]
  {examples/expected/m1_checkup.txt}

\section{How to read the numbers}

Each analysis section states its own criteria and caveats; the full criterion
reference is Chapter~\ref{ch:criteria}. Three points deserve emphasis here:

\begin{readbox}
\begin{itemize}
  \item \textbf{Silence is a result.} A clean output produces few findings. The
    findings you do get are graded: \sev{info} means ``perform this check'',
    \sev{warn} means a concrete doubt was found, \sev{error} means a conclusion
    must not be relied on, \sev{refuse} means the program declined to conclude.
  \item \textbf{The entropy spectrum in A2 is a bound, used one-sidedly.} A
    bound at or below the 0.14 line \emph{excludes} strong multireference
    character on that metric; a bound above it is \emph{indecisive}, not
    evidence for multireference character. The true four-state entropy needs the
    2-RDM route (Section~\ref{sec:menu1-local-spin}).
  \item \textbf{An empty result can be the point.} ``No f-composition orbital
    in the occupied partition'' on a supposed f system flags a wrong solution
    branch -- a known failure mode of f-block SCF/CASSCF (multiple solutions);
    change the initial guess or the active window and re-run.
\end{itemize}
\end{readbox}

\section{Worked example}

\subsection{A CASSCF output with warnings (N\textsubscript{2})}

The example above (\file{n2\_casscf\_nevpt2.out}, CASSCF(6,6) + SC-NEVPT2)
produces three warnings; Chapter~\ref{ch:report} shows them. Read them in
order: \sev{warn}~1 says the CASSCF reported convergence by the \emph{energy}
criterion only -- the orbital gradient was not checked, and the perturbation
layer depends on the orbitals; the action is to compare the printed gradient
against the tolerance, and tighten it if in doubt. Warnings 2 and 3 concern the
active space's near-empty and near-full orbitals (occupations 0.0020 and
1.9924): both are ``please confirm''-level -- for a full f shell or a magnetic
window this may be deliberate.

\subsection{An f\textsuperscript{1} system: the entropy bound is indecisive}

The same panel run on the generated Ce\textsuperscript{3+} input
(\file{generated\_ce3\_sarc2.out}) shows the open-shell limit of the bound.
Because the singly occupied 4f orbital has occupation 1, its bound is
$\ln 4 = 1.386294$ \emph{by construction} -- the true entropy of a single
unpaired electron is 0 or $\ln 2$ depending on the spin description, and the
bound cannot tell. The panel therefore reports the entropy indicator as
``above: indecisive'', lets the fractional-occupation count (0 of 7) carry the
evidence, and returns the combined verdict \emph{MR character not indicated}
with the T1 indicator named as unavailable.

\subsection{Local spin analysis}
\label{sec:menu1-local-spin}

When the input defined fragments (an atom label such as \code{N(1)}), ORCA
prints a \code{LOCAL SPIN ANALYSIS} block and menu 1 adds the A6 section:

\lstinputlisting[style=fbkcode, firstline=8, lastline=24,
  title={A6 section, block selection (from \file{n2\_stretch\_local\_spin.out})}]
  {examples/expected/products/n2_stretch_local_spin.out.fbk.md}

The full section reports the per-fragment $\langle S_z^A \rangle$ and
$S_{\mathrm{eff}}(A)$, the $\langle S_A S_B \rangle$ matrix, the total
$\langle S^2 \rangle$ with a spin-purity reading (this broken-symmetry example
gives 2.8938 instead of 0 -- intended non-eigenstate behaviour, not a file
defect), the metal-fragment spin share when an f-block fragment is present, and
an explicit boundary paragraph: \emph{local spin analysis is not the
environment spin polarisation entropy $\Delta S_E$}; the real $\Delta S_E$
needs the two-step off-line route (\code{orca\_2json} for the spin-resolved
2-RDMs plus \code{orca\_loc} localised orbitals), which the section spells out
together with the measured scale it stands in for.

\subsection{The diffuse-orbital (Rydberg) check}

Rule~G of the decision ladder asks whether the active orbitals are
Rydberg-contaminated, which an occupation-based selection cannot see.  The A8
section appears when the \emph{output itself} carries both the printed basis
(\code{!PrintBasis}) and the composition table; it then reports, per active
orbital, the diffusest and tightest Gaussian exponent of that orbital's
dominant (element, angular-momentum) channel, the orbital's weight on the
channel, and the ranking by diffuseness:

\lstinputlisting[style=fbkcode, firstline=68, lastline=88,
  title={A8 section (from the N\textsubscript{2} CAS(6,6) run reprinted with
  \code{!PrintBasis})}]
  {examples/expected/products/n2_diffuse.out.fbk.md}

The section deliberately carries no threshold: the source asks for
identification and sorting, and the exponent's absolute value does not transfer
between elements (measured: the diffusest s exponent is $0.19$ on N/def2-SVP
but $0.022$ on Eu/SARC2, whose valence functions are the diffuse end of \emph{that}
basis).  Read the diffusest exponent against the tightest one of the same
channel -- both are printed -- and across the active orbitals; the source's
remedy for a diffuse candidate is to extend the active space with it and
recompute the target states with SA-CASSCF + NEVPT2.

\subsection{A transition-state job}

Run on a frequency output, menu 1 adds the D4 findings -- one imaginary mode
reported as a first-order saddle (the transition-state candidate), together
with the reminder that a barrier derived from it must state its reference
point:

\lstinputlisting[style=fbkcode, firstline=13, lastline=18,
  title={Findings on \file{fhh\_optts\_freq.out} (excerpt)}]
  {examples/expected/products/fhh_optts_freq.out.fbk.md}

\section{Caveats, refusals and related sections}

\begin{refusalbox}
\begin{itemize}
  \item Menu 1 never changes your output file; it only writes the two report
    files beside it.
  \item It cannot check what the output does not contain: a DFT output has no
    CASSCF indicators, a CASSCF job prints no T1, and a job without declared
    fragments has no local spin block. The report names what is missing instead
    of guessing.
  \item Unrecognised files are rejected with a next-step message; parsing
    failures never produce an empty silent report.
\end{itemize}
\end{refusalbox}

Related: \menu{2} (geometry from a structure), \menu{7} (cross-level
consistency), \menu{9} (SCF rescue when the check-up reports SCF problems), and
Chapter~\ref{ch:criteria} for every threshold used above.
